\documentclass[12pt]{article}
\usepackage[utf8]{inputenc}
\usepackage[T1]{fontenc}
\usepackage[margin=1in]{geometry}
\usepackage{setspace}
\usepackage{xcolor}
\usepackage[colorlinks=true, urlcolor=blue!60!black, linkcolor=blue!60!black]{hyperref}
\setlength{\parindent}{0em}
\setlength{\parskip}{0.6em}
\renewcommand{\baselinestretch}{1.0}
\begin{document}
\begin{center}
{\LARGE \textit{Trump v. CASA, Inc.}, 606 U.S. 831 (2025)\par}
\vspace{0.5em}
{\large Case Brief --- Two-Page Edition\par}
\end{center}

\section{Memory Jogger}
Universal injunctions likely exceed federal courts' equitable power; relief is limited to what gives each plaintiff complete relief (\href{https://supreme.justia.com/cases/federal/us/606/24a884/}{Justia}).

\section{Detailed Case Facts}
Executive Order No.~14160 denied citizenship documents to many U.S.-born children of noncitizen parents. Individuals, organizations such as CASA, and 22 States sued in three districts under the Citizenship Clause and the Nationality Act.

\section{Procedural History}
Each district court entered a universal preliminary injunction; three circuits denied stays. The Government sought only partial stays, not defending the order.

\section{Judicial Votes}
Partial stays granted 6--3. Barrett wrote for the Court; Thomas, Alito, and Kavanaugh concurred; Sotomayor and Jackson dissented.

\section{Holding}
The Judiciary Act of 1789 does not likely authorize universal injunctions. The injunctions were stayed to the extent broader than necessary to give each plaintiff with standing complete relief; the order's legality was not decided.

\section{Analysis of Opinions}
\textbf{Majority:} founding-era equity gave party-specific relief; universal injunctions appeared only in the twentieth century; the bill of peace became the Rule 23 class action. Indivisible relief may incidentally help others, but one plaintiff's complete relief is not universal relief. \textbf{Concurrences:} complete relief is a ``ceiling'' (Thomas); beware state-standing and class loopholes (Alito). \textbf{Dissents:} the Government's ``gamesmanship'' leaves rights unprotected for nonparties (Sotomayor), creating a ``zone of lawlessness'' (Jackson).

\section{Examples: Future Applications}
\textbf{Same-Side:} two plaintiffs' nationwide injunction against a wire-transfer surcharge must be narrowed. \textbf{Opposite-Side:} a certified Rule 23(b)(2) class or APA set-aside remains available. \textbf{Fence-Sitter:} a State claiming that cross-border student migration requires broader relief.

\section{Critique}
Critics question freezing equity at 1789 practice, and note that classes, state suits, and APA vacatur have largely replicated nationwide relief.

\section{Key Quotations}
\begin{quote}
``\,`Complete relief' is not synonymous with `universal relief.' It is a narrower concept: The equitable tradition has long embraced the rule that courts generally `may administer complete relief between the parties.'\,''
\end{quote}
\noindent --- Justice Barrett, for the Court.

\end{document}
